% ------------------------------------------------------------------
%  Capitolo 17 — L'intelligenza artificiale: stampella o tutor?
%  Parte IV
%  Ricerca di riferimento: ricerca 01 §13; 02 §2.1, §6.4; 03 §2.3
%  Voci bibliografiche principali: bastani2025, kestin2025, kosmyna2025, risko2016
%  Epigrafe: ◐ verificata indirettamente (vedi ricerca 03 e registro, ricerca 00)
%  Scheda del capitolo: ricerca/06_indice-ragionato.md, capitolo 17
% ------------------------------------------------------------------
\chapter{L'intelligenza artificiale: stampella o tutor?}
\label{cap:ia}

\epigrafe[non hai trovato un farmaco per la memoria, ma per il richiamo]{οὔκουν μνήμης ἀλλὰ ὑπομνήσεως φάρμακον ηὗρες}{Platone, Fedro 275a}

\iniziale{Q}{uesto capitolo} \segnaposto{apertura: da scrivere — vedi la scheda nell'indice ragionato}.

\section{Il mito di Theuth}

\segnaposto{da scrivere}

\section{Che cosa dicono i primi esperimenti}

\segnaposto{da scrivere}

\section{Stampella o tutor}

\segnaposto{da scrivere}

\section{Regole per usare l'IA senza perdere l'apprendimento}

\segnaposto{da scrivere}

\begin{daricordare}
\segnaposto{il principio del capitolo in due righe}
\end{daricordare}

\begin{domande}
  \item \segnaposto{domanda di recupero su questo capitolo}
  \item \segnaposto{domanda di applicazione a una materia che studi}
  \item \segnaposto{domanda di ripresa su un capitolo precedente}
\end{domande}

\begin{sintesi}
  \item \segnaposto{punto di sintesi}
\end{sintesi}

\finecapitolo
